%Ficheiro Principal - Contem packages a usar e a definição do esqueleto do relatório.
% Recursos usados durante a elaboração deste documento:
% - https://www.overleaf.com/learn - Documentação do overleaf, serve de rampa de lançamento para aprender o LaTeX.
% - https://www.tablesgenerator.com/ - Permite a criação de tabelas para LaTeX sem ter que lidar com a sintaxe do mesmo.
% - https://app.diagrams.net/ - Permite criar figuras e diagramas de forma fácil. Permite a exportação das figuras criadas para PDF o que preserva a qualidade da mesma.
% - https://www.google.com/ - Para tudo o resto.

\documentclass[12pt,twoside,a4paper,openany]{book}
\usepackage[utf8]{inputenc}
\usepackage[a4paper,top=4.0cm,bottom=3.0cm,left=2.54cm,right=2.54cm,headheight=80pt,headsep=0.5cm,footskip=1.2cm]{geometry}

\usepackage[portuguese]{babel}
\usepackage{setspace}
\makeatletter
%APAGA A LINHA DO INÍCIO DAS PGs
\usepackage{ragged2e}
\usepackage{graphicx}
\usepackage{subcaption}
\usepackage[T1]{fontenc}
\usepackage{blindtext}
\usepackage[table,xcdraw]{xcolor}
\usepackage{soul} % para riscar texto com \st{}
\newcommand{\confidencial}{\textcolor{red}{[CONFIDENCIAL]}}

\usepackage{fancyhdr}
\usepackage{float} % para a imagem não saltar pgn
\usepackage{listings}
%\usepackage{minted}
\usepackage{csquotes}
\usepackage[open=true]{bookmark}
\usepackage{xcolor}
\usepackage{tikz} %para criar tarjas pretas em textos
\usetikzlibrary{positioning, shapes.geometric}
\usepackage{minitoc}


\lstset{
  language=Python,
  inputencoding=utf8,
  upquote=true, % Aspas retas
  columns=flexible, % Para a largura das letras ajustar bem
  extendedchars=true,
  basicstyle=\ttfamily\small,
  keywordstyle=\color{blue},
  commentstyle=\color{gray},
  stringstyle=\color{olive},
  numbers=left,
  frame=single,
  showstringspaces=false, 
  tabsize=4,
  breaklines=true,
  moredelim=**[is][\confidencial]{``}{``},
  literate={á}{{\'a}}1 {ã}{{\~a}}1 {â}{{\^a}}1 {à}{{\`a}}1
           {é}{{\'e}}1 {ê}{{\^e}}1 {í}{{\'i}}1 {ó}{{\'o}}1
           {ô}{{\^o}}1 {õ}{{\~o}}1 {ú}{{\'u}}1 {ç}{{\c{c}}}1
           {Ç}{{\c{C}}}1 {Á}{{\'A}}1 {É}{{\'E}}1 {Í}{{\'I}}1 
           {Ó}{{\'O}}1 {Ú}{{\'U}}1 {Â}{{\^A}}1 {Ê}{{\^E}}1 
           {Ô}{{\^O}}1 {Õ}{{\~O}}1 {Ç}{{\c{C}}}1,
}



\setcounter{minitocdepth}{2}
\mtcsettitle{minitoc}{Índice do Capítulo}
\setlength{\headheight}{80pt} %!!!!!!!!!!!!!!!!!!!
\bookmarksetup{
  numbered, 
}
%APAGA A LINHA DO INÍCIO DAS PGs
\usepackage{ragged2e}
\usepackage{titlesec}
\titleformat{\chapter}[hang]{\normalfont\huge\bfseries}{\thechapter.~}{0pt}{}
\titlespacing*{\chapter}{0pt}{-20pt}{20pt}
\titlespacing{\subsection}{0pt}{*1.5}{*0.8}

% with this we ensure that the chapter and section
% headings are in lowercase.
\renewcommand{\chaptermark}[1]{\markboth{#1}{}}
\renewcommand{\sectionmark}[1]{\markright{\thesection\ #1}}

\usepackage[style=apa]{biblatex}
\usepackage{tabularx}
\usepackage{pdfpages}
\usepackage{placeins}
\usepackage{svg}
\usepackage{amsmath}
\usepackage{color, colortbl}
\addbibresource{report/bibliografia.bib}
\usepackage{hyperref}
\def\blankpage{%
      \clearpage%
      \thispagestyle{empty}%
      %\addtocounter{page}{-1}% Rui - comentei para evitar abstract em página par
      \null%
      \clearpage}
\usepackage{caption}
%\DeclareCaptionType{equ}[Fórmula][]
\usepackage{acronym}


%PARA APAGAR A LINHA QUE APARECIA NO INÍCIO DE CADA PÁGINA
\renewcommand{\headrulewidth}{0pt} 


\usepackage{matlab-prettifier}
\usepackage[normalem]{ulem}

\newif\ifconfid
\confidtrue                   % comment this line out to unveil all hidden text
\colorlet{confidcolor}{black} % change the colour of boxes here (e.g. "red")


% ----------- implementation details -----------
\makeatletter

% from Martin's Scharrer's answer (https://tex.stackexchange.com/a/16004/21891)
\def\confid@UL@putbox{%
  \ifx\UL@start\@empty%
  \else % not inner
    \vrule\@width\z@ \LA@penalty\@M
    {\UL@skip\wd\UL@box \UL@leaders \kern-\UL@skip}%
    \phantom{\box\UL@box}%
  \fi
}

\newcommand\confidentiel{}%

\ifconfid
\renewcommand\confidentiel[1][confidcolor]{%
 \bgroup%
 \let\UL@putbox\confid@UL@putbox%
 \markoverwith{\hbox to.01em{\hss\textcolor{#1}{|}\hss}}\ULon%
}
\fi

\lst@AddToHook{Init}{%
  \ifconfid%
    \lstset{moredelim=[is][\confidentiel]``,keepspaces}%
  \else
    \lstset{moredelim=**[is][]``}%
  \fi
}

% Anteração no título do ToC
% https://tex.stackexchange.com/questions/28516/how-to-change-the-title-of-toc
\addto\captionsportuguese{% Replace "english" with the language you use
  \renewcommand{\contentsname}%
    {Índice}%
}

% PDF bookmarks  
% https://tex.stackexchange.com/questions/330534/how-to-properly-add-section-to-pdf-bookmarks
\makeatletter
\newcommand*{\sectionbookmark}[1][]{%
  \bookmark[%
    level=section,%
    dest=\@currentHref,%
    #1%
  ]%
}
\makeatother

%BEGIN: 2022-09-14 (Acerta a impressão fullduplex)
\let\tmp\oddsidemargin
\let\oddsidemargin\evensidemargin
\let\evensidemargin\tmp
\reversemarginpar
%END: 2022-09-14 (Acerta a impressão fullduplex)

% margens pela APA
%\usepackage[a4paper,margin=2.54cm]{geometry} 
\onehalfspacing % espaçamento entre linhas 1,5

\setlength{\footskip}{60pt}
\pagestyle{fancy}


\fancyhead[L]{
\includegraphics[height=2.4cm]{imagens/LOGO-ISEC.png}}
\fancyhead[LE,RO]{}
\fancyfoot[L]{\raisebox{2.2ex}{\tiny\parbox{10cm}{Mod6\_243.01\\Modelo C.27\_2.0\\SISTEMA INTERNO DE GARANTIA DE QUALIDADE}}}
\fancyfoot[C]{\hspace*{6cm}
\includegraphics[height=1.1cm]{imagens/barra_centro2030.png}} 
\fancyfoot[R]{\thepage}
%\setlength{\footskip}{25pt} % aumenta distância do rodapé

% CRIA: um estilo sem numeração só para capa e contracapa:
\fancypagestyle{semNumero}{
  \fancyhf{} % limpa tudo
  \fancyfoot[L]{\raisebox{2.2ex}{\tiny\parbox{10cm}{Mod6\_243.01\\Modelo C.27\_2.0\\SISTEMA INTERNO DE GARANTIA DE QUALIDADE}}}
  \fancyfoot[C]{\hspace*{6cm}
\includegraphics[height=1.1cm]{imagens/barra_centro2030.png}}
}

% CRIA: plain - footnote
\fancypagestyle{plain}{
  \fancyhf{} % limpa o cabeçalho e rodapé padrão do plain
  \fancyhead[L]{
\includegraphics[height=2.4cm]{imagens/LOGO-ISEC.png}}
  \fancyfoot[L]{\raisebox{2.2ex}{\tiny\parbox{10cm}{Mod6\_243.01\\Modelo C.27\_2.0\\SISTEMA INTERNO DE GARANTIA DE QUALIDADE}}}
  \fancyfoot[C]{\hspace*{6cm}
\includegraphics[height=1cm]{imagens/barra_centro2030.png}}
  \fancyfoot[R]{\thepage}
  \renewcommand{\headrulewidth}{0pt}
}


% Comando para "redactar" (tarja preta) do "\usepackage{tikz}"
% \redact{informações sensíveis}
\newcommand{\redact}[1]{%
  \tikz[baseline]{%
    \node[fill=black, text=black, inner sep=0.5ex, anchor=base] {#1};%
  }%
}


% Configuração da Lista de Figuras (LOF) e Lista de Tabelas (LOT)
\usepackage[titles]{tocloft}
\renewcommand{\cftfigpresnum}{Fig.~~}
\renewcommand{\cftfigaftersnum}{~-~}
\setlength{\cftfignumwidth}{4.5em}

\renewcommand{\cfttabpresnum}{Tabela~}
\renewcommand{\cfttabaftersnum}{~}
\setlength{\cfttabnumwidth}{5.0em}

\begin{document}
\frontmatter

%%%%%%%%%%%%%%%%%%%%%%%%%%%%%%%%%%%%%%%%%%%%%%%%%%%%%%%%%%%%%%%%%%%%%
%                            INÍCIO DA CAPA
\begin{center}

    

\includegraphics[width=0.4\textwidth]{imagens/LOGO-ISEC.png}\\

\vspace{1.5cm}
    
\Large\textbf{
Instituto Politécnico de Coimbra\\
\textit{Instituto Superior de Engenharia de Coimbra}}
\vspace{1.5cm}

\textbf{Curso Técnico Superior Profissional em\\ 
Tecnologias e Programação de Sistemas de Informação}
\vspace{1cm}

\textrm{Relatório de Estágio desenvolvido no âmbito da\\ Componente de Formação em Contexto de Trabalho}\\
\vspace{1cm}

\textbf{20 Recolher - Gestão de Resíduos,  Lda} \\   
\vspace{1cm}

\normalsize\textbf{Guilherme Laranjeiro Ventura}

\vspace{1cm}

\small\textrm{COIMBRA, julho de 2026}
\vspace{1cm}
   
\end{center}


\thispagestyle{semNumero} % Aplica o estilo
%\fancyfoot[CO]{\includegraphics[height=1cm]{imagens/Barra_PRR2.png}} 


%%%%%%%%%%%%%%%%%%%%%%%%%%%%%%%%%%%%%%%%%%%%%%%%%%%%%%%%%%%%%%%%%%%%%
%                           INÍCIO DA CONTRA CAPA

\thispagestyle{semNumero} % Aplica o estilo

\begin{center}
    
\includegraphics[width=0.4\textwidth]{imagens/LOGO-ISEC.png}\\

\vspace{1.5cm}
\Large\textbf{Instituto Politécnico de Coimbra\\
\Large\textit{Instituto Superior de Engenharia de Coimbra}}\\
\vspace{1.5cm}

\textrm{Relatório do Estágio desenvolvido no âmbito da\\ 
Componente de Formação em Contexto de Trabalho\\}

\small\textrm{04/03/2026 a 14/07/2026}

\vspace{1.5cm}

\textrm{Guilherme Laranjeiro Ventura\\
a2024149428@isec.pt\\
+351 935602163\\}

\end{center}

\vspace{1cm}
\begin{flushright}
\begin{minipage}{0.5\textwidth}
    \raggedright
\small\textbf{Orientador(es):}\\
\small\textrm{Paulo Miguel Gouveia Mariano\\
Professor, ISEC\\}
\small\textbf{Supervisor(es) na Empresa/Instituição:}\\
\small\textrm{Fernando Oliveira\\
CEO, 20 Recolher - Gestão de Resíduos,  Lda}
\end{minipage}
\end{flushright}




%%%%%%%%%%%%%%%%%%%%%%%%%%%%%%%%%%%%%%%%%%%%%%%%%%%%%%%%%%%%%%%%%%%%%
%                          AGRADECIMENTOS


%=============================================================================================

%=======================================================
%-----AGRADECIMENTOS
%=======================================================

\section*{Agradecimentos}
\sectionbookmark{Agradecimentos}

Gostaria de expressar o meu sincero agradecimento a todos os que, de forma direta ou indireta, contribuíram para a realização deste estágio e para a elaboração deste relatório.

À empresa 20 Recolher, em especial ao Diretor Executivo e supervisor de estágio, Fernando Oliveira, pela oportunidade concedida, pela confiança depositada nas minhas capacidades e pela disponibilidade demonstrada ao longo de todo o percurso.

Ao Instituto Superior de Engenharia de Coimbra (ISEC) e ao corpo docente do CTeSP de Tecnologias e Programação de Sistemas de Informação (TPSI), pelos ensinamentos e pela preparação técnica transmitida durante a formação académica.

Ao meu orientador de estágio no ISEC, pelo acompanhamento e apoio prestados no desenvolvimento deste documento.

Ao meu colega de estágio, pelo espírito de equipa, colaboração e entreajuda partilhados nas tarefas diárias.

À minha família e aos meus amigos, pelo apoio incondicional e incentivo demonstrados ao longo de todo este percurso.

%=============================================================================================


%=======================================================
%=======================================================



%%%%%%%%%%%%%%%%%%%%%%%%%%%%%%%%%%%%%%%%%%%%%%%%%%%%%%%%%%%%%%%%%%%%%
%                           ÍNDICE
\pagestyle{fancy} %plain - footnote


\pdfbookmark[section]{\contentsname}{toc}
%Indice
\dominitoc
%\addtocontents{toc}{\vspace{10pt}} % evita que fique colado ao rodapé
\tableofcontents

%%%%%%%%%%%%%%%%%%%%%%%%%%%%%%%%%%%%%%%%%%%%%%%%%%%%%%%%%%%%%%%%%%%%%
%                             LISTAS
%\addcontentsline{toc}{chapter}{Lista de Figuras}
\listoffigures
\sectionbookmark{Lista de Figuras}
\blankpage{}

%\addcontentsline{toc}{chapter}{Lista de Tabelas}
\listoftables
\sectionbookmark{Lista de Tabelas}

%%%%%%%%%%%%%%%%%%%%%%%%%%%%%%%%%%%%%%%%%%%%%%%%%%%%%%%%%%%%%%%%%%%%%
%                           ACRÓNIMOS

%Acronimos
%Nomenclatura
\chapter*{Nomenclatura}
\pdfbookmark[0]{Nomenclatura}{nomenclatura}
\fancyhead[LO]{\bfseries{NOMENCLATURA}}
\fancyhead[RE]{\bfseries{NOMENCLATURA}}
\section*{Abreviaturas}
\begin{acronym}

\acro{API}{\textit{Application Programming Interface}}
\acro{CCDRC}{Comissão de Coordenação e Desenvolvimento Regional do Centro}
\acro{CIM-RC}{Comunidade Intermunicipal da Região de Coimbra}
\acro{CTeSP}{Curso Técnico Superior Profissional}
\acro{ISEC}{Instituto Superior de Engenharia de Coimbra}
\acro{NVMe}{\textit{Non-Volatile Memory Express}}
\acro{RAM}{\textit{Random Access Memory}}
\acro{REEE}{Resíduos de Equipamentos Elétricos e Eletrónicos}
\acro{SATA}{\textit{Serial Advanced Technology Attachment}}
\acro{SEO}{\textit{Search Engine Optimization}}
\acro{SMART}{\textit{Self-Monitoring, Analysis and Reporting Technology}}
\acro{SSD}{\textit{Solid State Drive}}
\acro{TPSI}{Tecnologias e Programação de Sistemas de Informação}
\end{acronym}
\fancyhead{} % limpa cabeçalho
\renewcommand{\headrulewidth}{0pt} % remove linha horizontal
\blankpage{}


%%%%%%%%%%%%%%%%%%%%%%%%%%%%%%%%%%%%%%%%%%%%%%%%%%%%%%%%%%%%%%%%%%%%%
%                           CAPÍTULOS

\mainmatter %começa a numerar as páginas com números arábicos
\cleardoublepage
\chapter{Introdução}\label{chap1}
\acresetall


\section{Enquadramento e Motivação}
% Descreve o contexto geral do estágio, por que razão este projeto é importante
% e qual a motivação pessoal/profissional para a realização do mesmo.

Escreve aqui sobre o contexto geral do teu estágio.

\section{Apresentação da Entidade de Acolhimento}
% Apresenta a empresa "20 Recolher". O que faz, a sua área de atuação,
% onde se localiza e a sua importância no mercado. Podes adicionar secções
% sobre o organograma ou equipa se for relevante.

A 20 Recolher é uma empresa...

\section{Objetivos do Estágio}
% Define claramente os objetivos gerais e específicos que pretendes atingir
% ao longo do teu estágio na empresa.

Os principais objetivos definidos para este estágio incluem:
\begin{itemize}
    \item Objetivo 1;
    \item Objetivo 2;
\end{itemize}

\section{Estrutura do Relatório}
% Faz um resumo do conteúdo de cada capítulo do relatório para guiar o leitor.

O presente documento está estruturado da seguinte forma:
\begin{itemize}
    \item \textbf{Capítulo~\ref{chap1} -- Introdução:} Apresenta o enquadramento, a empresa de acolhimento, os objetivos e a organização do relatório;
    \item \textbf{Capítulo~\ref{chap2} -- Enquadramento Empresarial:} Apresenta a entidade de acolhimento (20 Recolher), a equipa de trabalho e o enquadramento do estagiário;
    \item \textbf{Capítulo~\ref{chap3} -- Atividades Planeadas:} Descreve os objetivos definidos e o plano de trabalhos inicial;
    \item \textbf{Capítulo~\ref{chap4} -- Atividades Desenvolvidas e Resultados:} Detalha as tarefas realizadas, dividindo-as entre o desenvolvimento de software (aplicações e base de dados) e o trabalho manual de hardware (testes de componentes e triagem de resíduos eletrónicos), apresentando os respetivos resultados obtidos;
    \item \textbf{Capítulo~\ref{chap5} -- Sugestões, Comentários e Propostas de Melhoria:} Reúne reflexões críticas, sugestões e propostas de evolução técnica para o projeto;
    \item \textbf{Capítulo~\ref{chap6} -- Conclusões:} Faz um balanço final dos objetivos atingidos e das competências adquiridas.
\end{itemize}

\cleardoublepage
\chapter{Enquadramento Empresarial}\label{chap2}


\section{Apresentação da Entidade de Acolhimento}
A 20 Recolher - Gestão de Resíduos, Lda. é uma empresa que trabalha na recolha e encaminhamento para reciclagem de \ac{REEE}, material informático e consumíveis de impressão (como tinteiros e toners).

A empresa nasceu a partir de uma ideia distinguida em abril de 2013 com o prémio de empreendedorismo da \ac{CIM-RC}. Com a licença da \ac{CCDRC}, começou a trabalhar em 2014 em Tentúgal. Em novembro de 2022, devido ao crescimento do volume de trabalho e à necessidade de instalações maiores, a empresa mudou-se para a Zona Industrial de Cantanhede (Núcleo 1, Lote 2, Fração E), onde opera atualmente.

O foco da 20 Recolher está no mercado empresarial e em entidades do Estado (como tribunais, escolas e conservatórias). O seu objetivo principal é garantir que os resíduos eletrónicos, que muitas vezes contêm substâncias perigosas para o ambiente, sejam recolhidos e tratados de forma correta, de forma a promover a reciclagem e a reutilização de componentes.

\begin{figure}[H]
    \centering
    
\includegraphics[width=0.4\textwidth]{imagens/logo-20-recolher.png}
    \caption{Logótipo da 20 Recolher - Gestão de Resíduos, Lda.}
    \label{fig:logo_empresa}
\end{figure}

\section{Equipa onde se insere o estagiário}
A equipa de desenvolvimento técnico da empresa é composta por dois estagiários: eu e um colega de curso, ambos alunos do \ac{CTeSP} de \ac{TPSI} no \ac{ISEC}. O trabalho decorre sob a supervisão de Fernando Oliveira, CEO e supervisor de estágio na 20 Recolher.

Por ser uma equipa pequena, a estrutura de trabalho é direta e informal. As decisões sobre a arquitetura do software, a escolha das tecnologias e a distribuição de tarefas foram tomadas de forma conjunta entre os dois estagiários, com a validação final do orientador de estágio.

\begin{figure}[H]
    \centering
    \begin{tikzpicture}[
        node distance = 1.2cm and 1cm,
        every node/.style = {draw, rectangle, rounded corners, minimum width=3cm, minimum height=1cm, align=center, fill=gray!10},
        line/.style = {draw, ->, thick}
    ]
        % Nodes
        \node (ceo) {\textbf{Fernando Oliveira} \\ \small Supervisor / CEO};
        \node (intern1) [below left=of ceo] {\textbf{Estagiário (Eu)} \\ \small Desenvolvimento Web};
        \node (intern2) [below right=of ceo] {\textbf{Colega de Estágio} \\ \small Desenvolvimento de Pesagens};
        
        % Connections
        \path [line] (ceo.south) -- ++(0,-0.5) -| (intern1.north);
        \path [line] (ceo.south) -- ++(0,-0.5) -| (intern2.north);
    \end{tikzpicture}
    \caption{Organograma da Equipa Técnica.}
    \label{fig:organograma}
\end{figure}

\section{Enquadramento Pessoal na Equipa}
Ao longo do estágio, as responsabilidades de desenvolvimento foram distribuídas de forma a responder às necessidades da 20 Recolher e a tirar partido do trabalho conjunto.

No início do estágio, trabalhei em colaboração com o meu colega no planeamento, levantamento de requisitos, desenho da arquitetura e modelação da base de dados da aplicação principal de pesagens e gestão de recolhas (desenvolvida em \texttt{.NET 10.0} e \texttt{Blazor Server}). Esta aplicação centraliza a gestão das recolhas e a integração com os terminais de pesagem instalados na zona de pesagens.

Posteriormente, para aumentar a produtividade e focar o desenvolvimento, dividimos o trabalho em duas partes:
\begin{itemize}
    \item O meu colega de estágio ficou responsável pelo desenvolvimento da aplicação de pesagens, que corre num ecrã tátil na zona de pesagens.
    \item Eu assumi a responsabilidade total pelo desenvolvimento e publicação do novo \textit{website} institucional da empresa, utilizando \texttt{Next.js 15} e \texttt{Tailwind CSS v4}.
\end{itemize}

Apesar desta divisão, mantivemos um contacto constante para garantir a integração dos dados entre as duas plataformas e a consistência visual. No desenvolvimento do \textit{website}, tomei de forma autónoma as decisões de design, estrutura e otimização de \ac{SEO}, com o objetivo de facilitar o pedido de recolhas online por parte de novos clientes.



\cleardoublepage
\chapter{Objetivos, Decisões Técnicas e Planeamento}\label{chap3}


O estágio curricular realizado na 20 Recolher --- Gestão de Resíduos, Lda. teve como ponto de partida um levantamento detalhado das necessidades tecnológicas da empresa.

\section{Estado Inicial e Problema Identificado}

No início do período de estágio, verificou-se que uma parte significativa dos processos operacionais dependia de procedimentos manuais e de informação distribuída por diferentes canais. Os pedidos de recolha eram, geralmente, comunicados através de telefone ou correio eletrónico, sem uma plataforma centralizada que permitisse a sua submissão, validação e acompanhamento estruturado.

O registo das pesagens era efetuado manualmente em papel, após a leitura visual dos valores apresentados pela balança industrial. Posteriormente, esse registo era encaminhado para a secretaria, onde a informação era tratada administrativamente. Este processo aumentava a possibilidade de erros de transcrição, perda ou duplicação de informação e dificultava a consulta do histórico das operações.

A presença digital da empresa também não disponibilizava uma área reservada através da qual os clientes pudessem submeter pedidos de recolha, acompanhar o respetivo estado ou consultar documentação associada de forma autónoma. Esta situação resultava numa maior dependência dos contactos telefónicos e das trocas de correio eletrónico.

Neste contexto, o estágio teve como principal finalidade a conceção e o desenvolvimento de um conjunto integrado de soluções tecnológicas destinado a centralizar a informação, apoiar os processos operacionais e melhorar a interação entre a empresa e os seus clientes.

\section{Objetivo Geral}

Conceber e desenvolver um ecossistema integrado de soluções de \textit{software} para apoiar e digitalizar os principais processos da 20 Recolher, desde a submissão de pedidos de recolha por parte dos clientes até ao registo das pesagens realizadas no armazém, promovendo a centralização da informação, a rastreabilidade das operações e a redução da dependência de procedimentos manuais.

\section{Objetivos Específicos}

Para concretizar o objetivo geral, definiram-se os seguintes objetivos específicos:

\begin{itemize}
    \item \textbf{Desenvolvimento do Website Institucional:} Criar um novo website institucional moderno para a empresa, com duas vertentes distintas: uma área pública para apresentação dos serviços e captação de novos clientes, e uma área reservada para clientes registados. Esta área reservada permite que os clientes efetuem pedidos de recolha de forma autónoma, acompanhem o estado das suas recolhas e comuniquem diretamente com os operadores através de um sistema de mensagens diretas. A presença digital estruturada é fundamental para reduzir o volume de contactos telefónicos e e-mails manuais, de forma a centralizar a gestão de clientes numa plataforma digital segura.
    \item \textbf{Desenvolvimento da Aplicação de \textit{Backoffice}:} Construir uma aplicação de gestão interna para os operadores da empresa, que centralize o controlo de clientes, a triagem e agendamento de pedidos de recolha, a visualização de pesagens e a administração de utilizadores e permissões. O objetivo é substituir o uso de folhas de cálculo e comunicações manuais.
    \item \textbf{Desenvolvimento do Terminal de Pesagem (\textit{Kiosk}):} Desenvolver uma aplicação destinada ao apoio das operações de pesagem e concebida para utilização através de um ecrã tátil. A aplicação permite ao operador introduzir manualmente os valores apresentados pela balança industrial e associá-los à recolha e ao respetivo tipo de material, reduzindo a necessidade de registos dispersos e facilitando a centralização da informação.

    \item \textbf{Modelação e Implementação da Base de Dados Centralizada:} Conceber e implementar um modelo de dados relacional capaz de suportar as diferentes aplicações do ecossistema, garantindo a integridade referencial e a centralização da informação. Sempre que aplicável, foi utilizada exclusão lógica, de forma a evitar a eliminação física imediata de registos relevantes e a preservar a sua disponibilidade para consulta.

    \item \textbf{Desenvolvimento da \ac{API} Central:} Criar uma \textit{Application Programming Interface} (API) \textit{RESTful} que funcione como a única camada de acesso à base de dados, garantindo que todas as aplicações comuniquem com os dados de forma centralizada, segura e consistente. Este objetivo surgiu da necessidade de separar as duas aplicações Blazor (backoffice e terminal de pesagem) que inicialmente foram concebidas como uma única aplicação.

    \item \textbf{Testes e Validação Funcional:} Validar o funcionamento das aplicações através de testes manuais realizados em computadores pessoais e num ecrã tátil disponibilizado pela empresa para futura utilização na zona operacional do armazém. Embora o equipamento ainda não se encontrasse instalado no local previsto, foi possível testar a interação por toque e simular o processo de introdução dos valores apresentados pela balança industrial. A integração direta com a balança não foi implementada, pelo que os valores continuaram a ser introduzidos manualmente pelo operador.
\end{itemize}


\section{Decisões Técnicas e Justificação das Escolhas}\label{sec:tecnologias-escolhas}

A escolha das tecnologias baseou-se em critérios como a adequação técnica ao problema, a facilidade de aprendizagem, a estabilidade e o suporte disponível no respetivo ecossistema, bem como a capacidade de integração entre as diferentes componentes do sistema.

\subsection{Desenvolvimento Frontend e Presença Web}
\begin{itemize}
    \item \textbf{\textit{Next.js} 16 (\textit{Website} Institucional):} O \textit{Next.js} foi escolhido por permitir desenvolver, no mesmo projeto, tanto as páginas institucionais como as áreas interativas do portal. Para isso, foram utilizados componentes de servidor e componentes de cliente, consoante as necessidades de cada funcionalidade. Os componentes de servidor são processados no servidor antes de o conteúdo ser enviado para o navegador, sendo adequados para apresentar informação estática, melhorar o desempenho e facilitar a indexação pelos motores de busca. Por sua vez, os componentes de cliente são executados no navegador do utilizador e permitem implementar funcionalidades interativas, como formulários, botões, filtros e atualizações dinâmicas da interface. A arquitetura baseada no \textit{App Router} permitiu ainda organizar o projeto de forma estruturada, através de rotas, páginas e componentes reutilizáveis.
    \item \textbf{\textit{Tailwind \ac{CSS}} v4 (Estilização):} O \textit{Tailwind CSS v4} foi adotado para acelerar o desenvolvimento visual da interface e garantir a sua adaptação a diferentes tamanhos de ecrã, incluindo smartphones, tablets e computadores. Esta tecnologia baseia-se em classes utilitárias, ou seja, pequenas classes previamente definidas que aplicam diretamente propriedades de estilo, como espaçamentos, cores, tamanhos ou alinhamentos. Desta forma, foi possível reduzir a necessidade de criar folhas de estilo extensas e personalizadas, diminuindo a quantidade de código CSS a manter e promovendo uma aparência consistente em todo o website. Como os clientes da empresa são maioritariamente entidades empresariais que podem consultar os seus dados através de vários dispositivos, a responsividade foi considerada um requisito operacional essencial e não apenas uma opção visual.
\end{itemize}

\subsection{Aplicações Empresariais Internas}
\begin{itemize}
    \item \textbf{\textit{Photino.Blazor} com .NET 10 (\textit{Backoffice} e \textit{Kiosk}):} O \textit{Photino.Blazor} foi utilizado no desenvolvimento das aplicações de ambiente de trabalho por permitir apresentar interfaces construídas com Blazor dentro de uma janela nativa do sistema operativo. Desta forma, foi possível desenvolver a interface com C\#, HTML e CSS, mantendo uma base tecnológica semelhante à utilizada na \ac{API} e aproveitando os conhecimentos previamente adquiridos em C\#. A aplicação recorre à \textit{WebView} disponibilizada pelo sistema operativo, isto é, a um componente que permite apresentar conteúdos web dentro de uma aplicação de ambiente de trabalho. Esta abordagem contribuiu para criar aplicações relativamente leves e adequadas aos equipamentos existentes na empresa.
    \item \textbf{\textit{Tauri} (Ferramenta de Toners):} O \textit{Tauri} foi utilizado para transformar a ferramenta de comparação de preços de toners numa aplicação de ambiente de trabalho destinada a uso interno. Esta tecnologia permite executar uma interface desenvolvida com tecnologias web dentro de uma janela nativa, recorrendo também à \textit{WebView} do sistema operativo. Esta escolha adequou-se ao objetivo de disponibilizar uma aplicação simples, leve e com reduzido consumo de recursos, que pudesse ser executada localmente no computador da empresa sem depender de um servidor web dedicado.
\end{itemize}

\subsection{Persistência e Infraestrutura de Dados}
\begin{itemize}
    \item \textbf{\textit{PostgreSQL} com \textit{Supabase}:} O PostgreSQL foi escolhido pela sua fiabilidade e pelo suporte a transações ACID, sigla correspondente a \textit{Atomicidade}, \textit{Consistência}, \textit{Isolamento} e \textit{Durabilidade}. Estas propriedades garantem que as operações são concluídas integralmente, respeitam as regras da base de dados, não interferem incorretamente entre si quando executadas em simultâneo e permanecem guardadas após a confirmação. Por exemplo, quando vários operadores registam pesagens ao mesmo tempo, o PostgreSQL controla a concorrência entre as transações, evitando registos incompletos ou alterações contraditórias. A plataforma \textit{Supabase} permitiu utilizar PostgreSQL gerido na nuvem, autenticação integrada e armazenamento de ficheiros, reduzindo a necessidade de configurar e manter servidores próprios.
    \item \textbf{Base de Dados Centralizada e \ac{API} Central:} O acesso aos dados foi concentrado numa \ac{API} desenvolvida em .NET 10. Esta camada recebe os pedidos das aplicações, valida os dados e as permissões dos utilizadores, aplica as regras de negócio e comunica com a base de dados. Desta forma, evitou-se a duplicação de lógica e melhorou-se o controlo e a manutenção do sistema.
    \item \textbf{\textit{Entity Framework Core} (\ac{ORM}):} O \textit{Entity Framework Core} foi utilizado como \ac{ORM}, sigla de \textit{Object-Relational Mapping}. Um ORM permite representar tabelas da base de dados através de classes em C\#, relacionando propriedades dos objetos com colunas e registos. Assim, foi possível consultar, adicionar e atualizar dados através de código C\#, sendo o \textit{Entity Framework Core} responsável por gerar os comandos SQL necessários. Foi utilizada uma abordagem \textit{Code-First}, em que o modelo da base de dados é definido no código e atualizado através de migrações, permitindo controlar e reproduzir as alterações do esquema nos diferentes ambientes do projeto.
\end{itemize}

\subsection{Fluxo de Trabalho e Publicação}
\begin{itemize}
    \item \textbf{\textit{Git} e \textit{GitHub}:} O controlo de versões com Git e o alojamento no GitHub foram adotados desde o primeiro dia de desenvolvimento. Esta decisão justifica-se pela necessidade de manter um histórico completo de todas as alterações ao código, possibilitar a reversão segura de mudanças problemáticas e garantir que o projeto não dependia de uma única máquina. Além disso, o Git permitiu que o desenvolvimento de novas funcionalidades ocorresse em ramos (\textit{branches}) isolados, sem perturbar a versão estável em produção.
    \item \textbf{\textit{Vercel} (Publicação Contínua do Website):} A plataforma Vercel foi escolhida para alojar o website institucional durante o desenvolvimento e os testes. A sua integração direta com o GitHub permitiu gerar automaticamente uma nova versão publicada após as atualizações enviadas para o repositório, tornando o website acessível sem necessidade de configurar ou manter servidores próprios. Esta abordagem facilitou a realização de testes em dispositivos reais, como computadores e smartphones, e permitiu verificar a responsividade e o funcionamento da interface ao longo do desenvolvimento.
    \item \textbf{Autenticação com \textit{Supabase Auth}:} O serviço \textit{Supabase Auth} foi utilizado para gerir o registo, a autenticação e as sessões dos utilizadores do portal. Esta opção permitiu recorrer a um serviço especializado, evitando o desenvolvimento integral de um mecanismo próprio de autenticação. O acesso às funcionalidades reservadas foi condicionado pela validação da sessão e das permissões associadas a cada utilizador.
\end{itemize}

\subsection{Validação Funcional e de Integração}
\begin{itemize}
    \item \textbf{Testes em Ambiente Real:} Validação da usabilidade e fluxo de trabalho do terminal quiosque tátil sob condições reais de armazém.
\end{itemize}


\section{Plano de Trabalhos}\label{sec:plano-trabalhos}

O plano de trabalhos foi estruturado em cinco fases, que se sobrepõem temporalmente de forma natural. As fases de análise e modelação ocorreram essencialmente no início do estágio, enquanto o desenvolvimento e os testes decorreram de forma contínua ao longo do período. O estágio decorreu entre 4 de março e 14 de julho de 2026.

\subsection{Fase 1: Levantamento de Requisitos e Análise}

Nesta fase inicial, procedeu-se à análise das necessidades operacionais da empresa através da observação direta dos processos existentes e das dificuldades identificadas no trabalho diário. O objetivo foi compreender o funcionamento da gestão de pedidos, do registo de pesagens e da comunicação com os clientes, identificando as necessidades que poderiam ser apoiadas através das soluções tecnológicas desenvolvidas.

Foram identificados e documentados requisitos de dois tipos:
\begin{itemize}
    \item \textbf{Requisitos Funcionais:} As funcionalidades que o sistema deveria disponibilizar, como a submissão de pedidos de recolha online, o registo manual assistido de pesagens, o sistema de mensagens entre clientes e operadores e a gestão de permissões por cargo;
    \item \textbf{Requisitos Não Funcionais:} As características de qualidade do sistema, como a segurança dos dados, a compatibilidade com dispositivos táteis, o tempo de resposta da \ac{API} e a disponibilidade numa infraestrutura alojada na nuvem.
\end{itemize}

\subsection{Fase 2: Desenho da Arquitetura e Modelação da Base de Dados}

Com base nos requisitos identificados, procedeu-se ao desenho da arquitetura global do sistema. A solução foi organizada em diferentes componentes, incluindo uma \ac{API} central, uma base de dados relacional, aplicações de ambiente de trabalho e um \textit{website}. A \ac{API} foi definida como a principal camada de comunicação com a base de dados, centralizando as operações de acesso e persistência da informação.

Paralelamente, foi desenhado o modelo de dados relacional, identificando as entidades principais, as suas relações de cardinalidade e as restrições de integridade referencial. O modelo foi implementado em \textit{PostgreSQL} através de uma abordagem \textit{Code-First} com \textit{Entity Framework Core}, sendo o esquema gerado e versionado através do mecanismo de migrações.

\subsection{Fase 3: Desenvolvimento e Implementação}

Esta foi a fase mais longa e correspondeu à implementação prática das soluções. O desenvolvimento foi realizado em entreajuda entre os dois estagiários da empresa, com a contínua revisão e alinhamento de decisões técnicas ao longo do processo. As componentes desenvolvidas abrangem:
\begin{itemize}
    \item \textbf{\textit{Website} Institucional e Portal de Cliente:} Desenvolvimento em \textit{Next.js} 16 e \textit{Tailwind \ac{CSS}} v4, com integração com a \ac{API} central;
    \item \textbf{Terminal de Pesagens (\textit{Kiosk}):} Desenvolvimento em \textit{Photino.Blazor} e .NET 10, com interface adaptada à utilização tátil e autenticação simplificada através de PIN;
    \item \textbf{\textit{Backoffice}:} Desenvolvimento da aplicação interna para gestão de clientes, pedidos, recolhas, pesagens e utilizadores;
    \item \textbf{\ac{API} Central:} Implementação das operações de acesso aos dados, autenticação, validação e comunicação entre as diferentes aplicações;
    \item \textbf{Ferramenta de Comparação de Toners:} Desenvolvimento de uma aplicação interna, posteriormente disponibilizada em formato de aplicação de ambiente de trabalho através do \textit{Tauri}.
\end{itemize}


\subsection{Fase 4: Testes e Validação}

Em paralelo com o desenvolvimento, foram realizados testes manuais e funcionais destinados a verificar os principais fluxos das aplicações e a comunicação entre as diferentes componentes do sistema.

A verificação abrangeu duas vertentes:
\begin{itemize}
    \item \textbf{Testes de Interface e Usabilidade:} Verificação da navegação, da introdução de dados e da adaptação da interface do terminal a um ecrã tátil disponibilizado pela empresa;
    \item \textbf{Testes de Integração de Dados:} Confirmação de que os dados submetidos através das diferentes aplicações eram corretamente enviados para a \ac{API} e persistidos na base de dados.
\end{itemize}

Estes testes permitiram identificar e corrigir problemas funcionais e de apresentação ao longo do processo de desenvolvimento. Não foram realizados testes formais com utilizadores finais, constituindo esta uma limitação do trabalho desenvolvido.


\subsection{Fase 5: Documentação e Entrega}

A fase final correspondeu à consolidação da documentação técnica e à preparação da entrega do projeto. Foram produzidos documentos de referência para futuros programadores, nomeadamente um guia de instalação em produção, um guia do \textit{frontend} para clientes e a documentação dos controladores da \ac{API}. Esta fase incluiu igualmente a elaboração do presente relatório de estágio.

\vspace{1em}
O cronograma das fases de trabalho ao longo do período de estágio encontra-se na Tabela~\ref{tab:cronograma}.

\begin{table}[H]
    \centering
    \caption{Cronograma das Fases do Plano de Trabalhos (2026).}
    \label{tab:cronograma}
    \begin{tabularx}{\textwidth}{|X|c|c|c|c|c|}
        \hline
        \rowcolor{gray!20} \textbf{Fase / Atividade} & \textbf{Mar.} & \textbf{Abr.} & \textbf{Mai.} & \textbf{Jun.} & \textbf{Jul.} \\ \hline
        Fase 1: Levantamento de Requisitos e Análise & \cellcolor{gray!50} X & & & & \\ \hline
        Fase 2: Arquitetura e Modelação da Base de Dados & \cellcolor{gray!50} X & \cellcolor{gray!50} X & & & \\ \hline
        Fase 3: Desenvolvimento e Implementação & & \cellcolor{gray!50} X & \cellcolor{gray!50} X & \cellcolor{gray!50} X & \cellcolor{gray!50} X \\ \hline
        Fase 4: Testes e Validação & & \cellcolor{gray!50} X & \cellcolor{gray!50} X & \cellcolor{gray!50} X & \cellcolor{gray!50} X \\ \hline
        Fase 5: Documentação e Entrega & & & & \cellcolor{gray!50} X & \cellcolor{gray!50} X \\ \hline
    \end{tabularx}
\end{table}
\cleardoublepage
\chapter{Atividades Desenvolvidas e Resultados}\label{chap4}
\thispagestyle{fancy}

Neste capítulo, são descritas em pormenor as atividades realizadas ao longo do período de estágio, as quais se dividiram em duas vertentes principais: a componente de engenharia e desenvolvimento de software (que engloba a modelação de base de dados, desenvolvimento do website e de aplicações empresariais) e a componente prática de hardware (que compreende tarefas manuais de triagem, teste e reciclagem de componentes eletrónicos nas instalações da empresa).

\section{Desenvolvimento de Software e Aplicações}

Esta secção descreve a conceção, arquitetura e implementação das soluções de software que compõem o ecossistema tecnológico desenvolvido para a empresa.

\subsection{Arquitetura Geral do Sistema}

O ecossistema tecnológico foi desenhado sob uma arquitetura distribuída e modular, visando garantir a separação de conceitos, a escalabilidade das aplicações e a facilidade de manutenção a longo prazo. O núcleo da solução assenta num modelo de \textit{Clean Architecture} (ou Arquitetura Cebola) implementado no backoffice, complementado por uma API centralizada que funciona como o único ponto de contacto com a base de dados.

A infraestrutura é constituída por três camadas lógicas principais:
\begin{itemize}
    \item \textbf{Camada de Apresentação (Frontend/Clientes):} constituída pelo Website Institucional em Next.js 15, pela aplicação de Backoffice e pelo Terminal de Pesagem (Kiosk), ambos desenvolvidos em Blazor Server no ecossistema .NET 10.
    \item \textbf{Camada de Serviços (API Central):} uma API RESTful em .NET 10 que expõe os \textit{endpoints} de negócio necessários e assegura que todas as validações, regras e acessos à base de dados ocorrem sob uma mesma lógica corporativa.
    \item \textbf{Camada de Dados e Infraestrutura (Supabase/PostgreSQL):} alojamento na nuvem da base de dados PostgreSQL e do serviço de armazenamento de ficheiros do Supabase.
\end{itemize}

O isolamento das regras de negócio na API central garante que as aplicações clientes (Web, Backoffice e Quiosque) não comuniquem diretamente com a base de dados, reduzindo o risco de acessos não autorizados e facilitando futuras integrações com outros sistemas da empresa.

\subsection{Fluxos Principais da Aplicação}

Para assegurar a digitalização eficiente dos processos da empresa, foram modelados e implementados três fluxos fundamentais:
\begin{enumerate}
    \item \textbf{Fluxo de Submissão e Triagem de Pedidos de Recolha:} O cliente acede à área reservada do website e submete um formulário indicando o material a recolher, anexando fotografias. O pedido entra na base de dados com o estado \textit{Pendente}. O operador, na app de Backoffice, recebe o alerta, analisa os dados e as fotos e agenda a recolha física, alterando o estado para \textit{Agendada}.
    \item \textbf{Fluxo de Pesagem e Integração Física:} No armazém, o motorista descarrega os resíduos. O operador no terminal de pesagem (Kiosk) escolhe a recolha ativa. O operador insere o valor do peso medido no terminal, que submete os dados à API central. A pesagem é associada ao pedido de recolha original, atualizando a conta corrente do cliente.
    \item \textbf{Fluxo de Alteração Cadastral e Auditoria de Segurança:} Visando impedir fraudes e dados cadastrais incorretos na faturação, sempre que o cliente edita os seus dados fiscais ou morada na área reservada do website, a atualização não é imediata. A alteração entra numa tabela de propostas temporárias. O operador de backoffice analisa o pedido e valida. Se aprovar, a API executa a persistência na tabela central de Clientes e cria um registo de auditoria (\textit{audit log}).
\end{enumerate}

\subsection{Modelação da Base de Dados}

A base de dados constitui o núcleo central de toda a infraestrutura tecnológica desenvolvida, sustentando de forma unificada tanto o website institucional (através da área reservada de cliente e do painel de gestão editorial) como as aplicações empresariais de backoffice e do terminal de pesagem. Para responder aos requisitos de consistência, integridade de dados e escalabilidade, optou-se pela utilização de um sistema de gestão de bases de dados relacional, especificamente o PostgreSQL. A base de dados é gerida e alojada na nuvem através da plataforma Supabase, o que confere elevada disponibilidade e simplifica as operações de monitorização e cópias de segurança.

A modelação física e a comunicação com o motor de base de dados foram estruturadas através de uma abordagem \textit{Code-First} suportada pelo Entity Framework Core (EF Core), o mapeador objeto-relacional (ORM) padrão da plataforma .NET. Sob esta metodologia, todo o esquema de tabelas, chaves primárias, chaves estrangeiras e restrições de integridade foi modelado sob a forma de classes de domínio na linguagem C\#. As atualizações estruturais e a sincronização do esquema com o servidor PostgreSQL foram geridas de forma estrita através do sistema de migrações (\textit{migrations}) do EF Core, assegurando a rastreabilidade e a reprodutibilidade do esquema em ambientes de desenvolvimento e produção.

A arquitetura de dados divide-se nas seguintes entidades principais:
\begin{itemize}
    \item \textbf{Cliente:} Armazena os dados de contacto, localização e de identificação necessários para a faturação e conformidade legal das entidades clientes. Esta entidade está associada ao utilizador registado no serviço de autenticação central do Supabase através de um identificador único de utilizador.
    \item \textbf{CodigoLer e Material:} Representam a catalogação regulamentar e comercial de resíduos. A tabela \textit{CodigoLer} armazena os códigos oficiais estabelecidos na Lista Europeia de Resíduos (LER) e a respetiva descrição. A tabela \textit{Material} associa um nome operacional e a informação tarifária aplicável a cada tipo de resíduo sob um código LER.
    \item \textbf{PedidoRecolha:} Regista as solicitações de recolha submetidas pelos clientes através do website ou inseridas manualmente pelos operadores. Guarda o estado de processamento do pedido, notas adicionais, informações de contacto e dados geográficos da recolha. Associa-se aos materiais pretendidos através da tabela de junção \textit{PedidoMaterialPretendido}.
    \item \textbf{Recolha:} Modela o evento logístico de recolha de resíduos. Contém a data programada, o estado de execução, dados geográficos e de contacto para a operação. Relaciona-se com os materiais selecionados para a recolha através da tabela de junção \textit{RecolhaMaterialPreSelecionado}.
    \item \textbf{Pesagem:} Regista as pesagens efetuadas na balança física. Guarda a quantidade líquida pesada, o valor calculado da transação, o material correspondente, o operador que validou a pesagem e a recolha associada.
    \item \textbf{Utilizador e Cargo:} Representam os utilizadores internos do sistema (como operadores de balança, funcionários de armazém e administradores). Cada conta está associada a um cargo, cujas permissões de acesso a funcionalidades do sistema são configuradas de forma granular através de um campo estruturado em formato JSON.
    \item \textbf{KioskDevice e KioskPairingCode:} Suportam a segurança dos quiosques físicos de pesagem. Registam os identificadores exclusivos dos terminais autorizados e gerem os códigos temporários utilizados para emparelhar em segurança novos dispositivos quiosque à infraestrutura central.
    \item \textbf{Noticias e RevisoeseNoticias:} Suportam o sistema de publicação de conteúdos do website institucional, gerindo dados textuais, categorias, referências a imagens de capa e campos destinados à otimização em motores de busca (SEO). A tabela \textit{RevisoeseNoticias} mantém o histórico de modificações para auditoria e reversão.
    \item \textbf{ChatMessage:} Implementa o suporte direto e bidirecional de mensagens entre os clientes e os operadores de backoffice, possuindo um mecanismo de expiração de dados com base na data de validade do registo.
    \item \textbf{SolicitacaoAlteracaoPerfil:} Modela o fluxo de controlo de alterações cadastrais iniciadas por clientes na área reservada. As propostas de alteração de dados de identificação e faturação são registadas temporariamente nesta tabela, exigindo a aprovação explícita de um operador no backoffice antes de atualizar a ficha do cliente.
    \item \textbf{Anexo:} Tabela polimórfica que armazena referências a ficheiros físicos (como fotografias de resíduos pesados ou comprovativos em formato PDF). Utiliza referências estruturadas para associar dinamicamente os ficheiros às entidades de negócio correspondentes.
    \item \textbf{AuthEmailQueue:} Fila assíncrona que processa o envio de e-mails do sistema (como convites de registo ou códigos de validação de adesão), gerindo a resiliência no envio com políticas de novas tentativas em caso de erro.
\end{itemize}

O modelo de dados implementado reflete ainda diversas decisões de design concebidas para otimizar o desempenho, a segurança e a utilidade operacional do sistema:
\begin{itemize}
    \item \textbf{Mecanismo de Exclusão Lógica:} De forma a manter a integridade dos dados históricos da atividade da empresa e cumprir requisitos de auditoria, nenhuma informação de negócio é eliminada fisicamente da base de dados. Todas as tabelas integram colunas destinadas a identificar o estado de eliminação lógica do registo e a autoria dessa ação. Quando um operador remove um registo, este é apenas marcado como inativo, ficando oculto das listagens regulares mas disponível para consulta histórica e recuperação no painel de administração.
    \item \textbf{Numeração Sequencial Relativa ao Cliente:} Para simplificar a comunicação de encomendas, implementou-se um número sequencial relativo a cada cliente nas entidades de pedidos e recolhas. Em vez de expor os identificadores numéricos globais gerados pela base de dados, cada cliente visualiza a sua atividade numerada consecutivamente (por exemplo, Pedido \#1 e Pedido \#2 do Cliente X), isolando a lógica sequencial de cada conta de cliente.
    \item \textbf{Gestão de Ficheiros Descentralizada:} Os anexos e imagens de capa de notícias são guardados em baldes (\textit{buckets}) de armazenamento de objetos do Supabase (compatível com S3). O registo do anexo na base de dados armazena apenas a referência única de acesso e metadados. Adicionalmente, a aplicação assegura que os ficheiros de imagem sofrem processos de compressão e conversão para o formato WebP antes de serem enviados, otimizando o consumo de tráfego de rede e armazenamento na nuvem.
    \item \textbf{Índices de Performance Direcionados:} Foram desenhados índices compostos personalizados para acelerar consultas cruciais. Na tabela de pesagens, por exemplo, desenhou-se um índice composto pelas colunas frequentemente consultadas em conjunto, como datas, identificadores de materiais e o estado lógico do registo, maximizando a eficiência de processamento na geração de relatórios e análises de desempenho.
\end{itemize}

O diagrama relacional simplificado da base de dados, detalhando as tabelas e as respetivas relações de cardinalidade, encontra-se representado na Figura~\ref{fig:erd}.

\begin{figure}[H]
    \centering
    \includegraphics[width=0.9\textwidth]{imagens/erd.png}
    \caption{Diagrama de Entidade-Associação da base de dados centralizada.}
    \label{fig:erd}
\end{figure}

\subsection{Website Institucional}

O website institucional foi desenvolvido utilizando a versão 15 do Next.js com a arquitetura App Router. Esta escolha tecnológica permitiu tirar partido de renderização do lado do servidor (\textit{Server-Side Rendering}) e geração estática de páginas para otimizar o tempo de resposta inicial, garantindo simultaneamente boas práticas de indexação em motores de busca (SEO). A estilização e o desenho responsivo da interface foram concebidos com a versão 4 do Tailwind CSS, resultando numa interface moderna, fluida e adaptável a dispositivos móveis, tablets e computadores.

A plataforma está organizada em três grandes áreas funcionais e de acesso regulado: a área pública de informação e captação de clientes, a área reservada de clientes e o painel de administração de conteúdos.

\subsubsection{Área Pública}
A secção pública do website destina-se a apresentar os serviços da empresa e a estabelecer canais de comunicação com a comunidade e potenciais parceiros de negócio, sendo composta pelas seguintes páginas:
\begin{itemize}
    \item \textbf{Página Inicial:} Apresenta uma secção introdutória de forte impacto visual, um resumo dos serviços principais de recolha de resíduos industriais, estatísticas gerais sobre o impacto ecológico das atividades da empresa e uma secção de perguntas frequentes (FAQs) para responder rapidamente a dúvidas recorrentes dos utilizadores.
    \item \textbf{Serviços e Institucional:} Páginas dedicadas a detalhar os procedimentos regulamentares da empresa, os tipos de resíduos recolhidos e a história da organização.
    \item \textbf{Notícias:} Apresenta uma listagem cronológica de artigos publicados pela empresa, dotada de paginação e filtros de pesquisa por categorias, com redirecionamento para artigos de leitura individual.
    \item \textbf{Contactos:} Disponibiliza um canal direto de envio de mensagens com validação de dados e a integração de um mapa interativo para localização física das instalações da empresa.
    \item \textbf{Política de Privacidade:} Documento regulamentar para assegurar a conformidade da recolha e tratamento de dados dos utilizadores face à legislação em vigor.
\end{itemize}

\subsubsection{Área Reservada de Cliente}
Acedida através de mecanismos de autenticação segura, a área de cliente faculta aos parceiros comerciais ferramentas para gerir autonomamente as suas operações diárias com a empresa, dividindo-se nos seguintes módulos funcionais:
\begin{itemize}
    \item \textbf{Painel de Controlo:} Mostra um resumo dinâmico das recolhas efetuadas, pedidos ativos e alertas de suporte pendentes.
    \item \textbf{Pedidos de Recolha:} Permite a submissão de novas solicitações de recolha mediante um formulário onde o cliente indica os materiais de resíduos que deseja recolher, dados de contacto adicionais e anexa documentos ou fotografias relevantes para a operação. Disponibiliza ainda a consulta, cancelamento e reativação de pedidos em curso.
    \item \textbf{Histórico de Recolhas:} Regista os eventos logísticos concluídos e agendados. Apresenta o peso líquido total registado na balança física e a respetiva codificação legal (Lista Europeia de Resíduos), permitindo a exportação de dados em formatos estruturados como PDF e tabelas CSV.
    \item \textbf{Catálogo de Materiais:} Listagem paginada e pesquisável de todos os resíduos autorizados e geridos pela empresa, permitindo a exportação rápida dos dados tarifários vigentes para formato CSV.
    \item \textbf{Perfil de Cliente:} Permite ao parceiro solicitar atualizações aos seus dados cadastrais e fiscais. Estas alterações são submetidas a um fluxo de controlo interno, exigindo a aprovação manual de um operador no backoffice antes de serem efetivamente integradas na base de dados, prevenindo dados incorretos e tentativas de fraude.
    \item \textbf{Suporte Técnico (Chat):} Canal de mensagens direto e em tempo real com a equipa da empresa, concebido para acelerar a resolução de litígios ou o esclarecimento de dúvidas sobre agendamentos de recolha.
\end{itemize}

\subsubsection{Painel de Administração}
O painel administrativo é uma secção restrita e protegida por perfis de acesso elevados (administrador e super-administrador do website), que permite a gestão direta dos conteúdos editoriais do website:
\begin{itemize}
    \item \textbf{Módulo de Notícias:} Permite a criação, edição e eliminação lógica de publicações utilizando um editor visual de texto estruturado. Inclui campos dedicados para a configuração de títulos, slugs de URL e meta-descrições para a otimização em motores de busca (SEO).
    \item \textbf{Controlo e RevisoeseNoticias:} Possibilita a visualização do histórico de revisões dos artigos, permitindo comparar alterações estruturais e reverter edições antes da sua publicação definitiva no website.
\end{itemize}

\subsubsection{Alojamento e Publicação}
Durante as fases de desenvolvimento e testes de validação com a equipa da empresa, o website institucional foi alojado na plataforma Vercel, beneficiando de processos automatizados de compilação e publicação contínua (\textit{Continuous Deployment}). A arquitetura foi concebida para permitir a transição simplificada para o domínio e servidores corporativos permanentes da empresa.

\subsection{Ferramenta de Comparação de Preços de Toners}

No âmbito das necessidades de otimização do processo de vendas de consumíveis da empresa, foi desenvolvida uma ferramenta utilitária dedicada à comparação de preços de toners. Esta ferramenta permite carregar tabelas tarifárias e propostas de compra de múltiplos potenciais clientes corporativos e parceiros, frequentemente fornecidas em formatos de folhas de cálculo com diferentes estruturas de colunas, e efetuar o mapeamento flexível e a normalização de referências de produto. Com isto, a 20 Recolher consegue cruzar a informação do seu inventário de toners com a procura do mercado e com as tabelas de preços da concorrência, identificando a proposta de venda comercialmente mais vantajosa e competitiva para os seus compradores.

A aplicação foi construída utilizando a biblioteca React e a ferramenta de compilação Vite para a interface web. O processamento dos ficheiros locais é assegurado pela biblioteca XLSX, enquanto o armazenamento do histórico de tendências de preços é gerido localmente no navegador através de IndexedDB. Adicionalmente, foi implementada uma camada de empacotamento nativo com o Tauri para possibilitar a sua distribuição e execução como aplicação de ambiente de trabalho (\textit{desktop}). O desenvolvimento desta ferramenta foi realizado em parceria com o colega de estágio, tendo os testes unitários do motor de normalização de dados sido implementados com a biblioteca Vitest.

\subsection{App de Backoffice}

Para além do portal de clientes, o ecossistema tecnológico integra a aplicação de gestão interna (Backoffice), desenvolvida em .NET 10 sob a tecnologia Blazor Server. Esta aplicação destina-se ao controlo administrativo global da atividade e é utilizada pelos operadores e administradores da empresa para gerir e monitorizar todo o fluxo logístico e de pesagens.

A aplicação encontra-se estruturada nos seguintes módulos funcionais principais:
\begin{itemize}
    \item \textbf{Gestão de Entidades:} Permite o controlo de registos cadastrais de clientes e utilizadores do sistema, bem como a definição e atribuição de cargos e permissões de acesso estruturadas de forma granular para garantir a segurança lógica do sistema.
    \item \textbf{Controlo Logístico:} Módulo dedicado à triagem e agendamento de pedidos de recolha de resíduos, planeamento de recolhas e visualização dos dados consolidados de pesagens e quantidades processadas.
    \item \textbf{Auditoria de Segurança:} Regista e expõe os diários de atividade (\textit{audit logs}), detalhando todas as operações estruturais realizadas na base de dados (como criação, modificação ou inativação lógica de registos), identificando o utilizador responsável e o método de autenticação utilizado.
    \item \textbf{Validação de Alterações Cadastrais:} Canal onde os operadores analisam e validam ou rejeitam as solicitações de alteração de dados de identificação e faturação submetidas pelos clientes através da área reservada do website.
\end{itemize}

Em termos de usabilidade, a interface foi desenhada para se adaptar a diferentes cenários de trabalho, incluindo a utilização em dispositivos móveis e tablets no terreno. Para tal, implementaram-se tabelas de dados dinâmicas com colunas configuráveis e persistidas localmente para cada operador, além de menus de ações acionados por toque longo (\textit{long press}), adequados a ecrãs táteis. A aplicação disponibiliza ainda a exportação de relatórios de atividade em formato PDF e um painel dedicado à recuperação de registos inativos através do mecanismo de exclusão lógica (\textit{Soft-Delete}). O desenvolvimento desta aplicação de gestão foi realizado em colaboração com o colega de estágio.

\subsection{App de Balança - Terminal Kiosk}

Para além do backoffice administrativo, o sistema dispõe de um módulo dedicado especificamente à zona física de pesagens, designado por Terminal Kiosk. Trata-se de uma aplicação desenvolvida em .NET 10 sob a tecnologia Blazor Server, com uma interface otimizada para interação tátil (\textit{touchscreen}), concebida para ser operada de forma rápida e intuitiva por funcionários no terreno.

As características estruturais e operacionais deste terminal incluem:
\begin{itemize}
    \item \textbf{Integração com Balança Física:} O sistema comunica em tempo real com a balança industrial da empresa através de portas de comunicação série (protocolos RS-232), efetuando a leitura e importação automática do peso físico sem necessidade de introdução manual por parte do operador, o que minimiza erros humanos e otimiza o fluxo de registo.
    \item \textbf{Autenticação Simplificada:} A autenticação no quiosque é baseada na introdução de um código numérico pessoal (PIN) através de um teclado virtual sobredimensionado no ecrã. Isto evita a necessidade de introduzir palavras-passe complexas em ecrãs táteis expostos a poeira e condições operacionais de armazém.
    \item \textbf{Processo Guiado de Pesagem (Wizard):} O registo de pesagens é estruturado na forma de um assistente passo-a-passo (\textit{Wizard}). O operador seleciona a recolha em curso, o material que está a ser pesado, confirma a leitura automática do peso e submete o registo de forma sequencial e validada estruturalmente pelas regras de negócio.
    \item \textbf{Emparelhamento de Terminais (Pairing Code):} Para garantir que apenas quiosques físicos autorizados comunicam com a base de dados central, implementou-se um fluxo de emparelhamento obrigatório. O administrador gera um código de autorização temporário no backoffice e introduz este código no novo terminal para estabelecer a ligação encriptada definitiva.
\end{itemize}

A aplicação partilha componentes reutilizáveis através do módulo comum do projeto e foi desenhada com foco na usabilidade e fiabilidade das operações sob condições de armazém. O desenvolvimento deste módulo foi realizado em estreita parceria com o colega de estágio.

\subsection{Integração entre Sistemas e API Central}

Inicialmente, as aplicações de backoffice administrativo e o terminal de pesagem física (Kiosk) foram concebidos e implementados de forma acoplada numa única aplicação. No entanto, a evolução do projeto e a necessidade de isolar os ambientes físico e de escritório exigiram a separação destas ferramentas em duas aplicações independentes. Para suportar esta nova topologia sem comprometer a consistência dos dados, foi desenvolvida a API Central (\texttt{RecolherApp.Api}) na plataforma .NET.

Esta API funciona como a única interface de comunicação com a base de dados relacional e expõe um conjunto de serviços RESTful consumidos por todos os componentes do ecossistema, incluindo o website institucional. A arquitetura global do sistema segue os princípios da \textit{Clean Architecture}, dividindo-se em seis camadas estruturais:
\begin{itemize}
    \item \textbf{\texttt{RecolherApp.Api}:} Camada de exposição que gere os pontos de entrada (\textit{endpoints}), tratamento de pedidos HTTP e serialização de dados.
    \item \textbf{\texttt{RecolherApp.Core}:} Núcleo do domínio que contém a definição das entidades de negócio, validações estruturais e interfaces de repositórios.
    \item \textbf{\texttt{RecolherApp.Infrastructure}:} Implementação de acessos estruturais a bases de dados (via EF Core) e integrações com o serviço de armazenamento de objetos (Supabase Storage).
    \item \textbf{\texttt{RecolherApp.UI}:} Interface de utilizador do painel de backoffice (Blazor Server).
    \item \textbf{\texttt{RecolherApp.UI.Pesagens}:} Interface de utilizador do quiosque de balança (Blazor Server).
    \item \textbf{\texttt{RecolherApp.UI.Shared}:} Biblioteca partilhada contendo componentes de interface reutilizáveis (como assistentes, menus e inputs táteis) e tipos de dados comuns às duas aplicações Blazor.
\end{itemize}

O website institucional em Next.js também comunica de forma estrita com esta API central, sem nunca aceder diretamente à base de dados PostgreSQL. O fluxo de integração de pedidos de recolha funciona da seguinte forma:
\begin{enumerate}
    \item O cliente preenche e submete o formulário de recolha na área reservada do website.
    \item O servidor Next.js recebe o pedido e utiliza uma camada de proxy interna (\texttt{src/proxy.js}) para reencaminhar a solicitação para a API .NET, anexando de forma segura o token JWT de autenticação do cliente.
    \item A API central valida os dados estruturais do pedido, atualiza a base de dados relacional e aciona a fila de processamento assíncrono para o envio automático de e-mails cadastrais.
    \item O novo pedido passa imediatamente a estar visível no painel administrativo de backoffice para fins de agendamento logístico.
\end{enumerate}

Adicionalmente, a API disponibiliza o catálogo de materiais regulamentares para popular as seleções de formulários no website, integrando um mecanismo de salvaguarda local (\textit{local fallback}) para manter o formulário de pedidos funcional em caso de instabilidade temporária na ligação de rede. O envio de ficheiros e comprovativos no website é reencaminhado pelo proxy para a API .NET, que executa tarefas automáticas de compressão e conversão de imagem para o formato WebP antes de submeter os ficheiros para o armazenamento permanente no Supabase Storage.

\subsection{Segurança e Autenticação}

Para proteger os dados cadastrais, comerciais e operacionais geridos pelas diferentes aplicações, foi implementado um modelo de segurança e autenticação estruturado em duas vertentes principais: uma destinada ao portal web público e de clientes, e outra direcionada às ferramentas internas da empresa (Backoffice e quiosque de pesagens).

\subsubsection{Segurança no Website e Área de Cliente}
A autenticação e autorização na área de cliente e no painel administrativo do website são geridas através da integração com o serviço centralizado Supabase Auth, estruturando-se da seguinte forma:
\begin{itemize}
    \item \textbf{Criação de Conta via Convite:} O processo de registo de um novo cliente é acionado automaticamente aquando da submissão do seu primeiro pedido de recolha de resíduos. O sistema regista os dados na base de dados e insere uma tarefa na fila de envio de e-mails, disparando um convite seguro contendo um link de ativação único. O utilizador utiliza esse link para definir a sua palavra-passe e ativar a sua área reservada.
    \item \textbf{Gestão de Sessões com Tokens JWT:} Após a autenticação com sucesso, o sistema emite tokens JSON Web Token (JWT) que atestam a identidade e perfil do utilizador. Estes tokens são transmitidos e armazenados de forma segura utilizando cookies configurados com a propriedade \textit{HttpOnly}, garantindo que o código JavaScript do lado do cliente não tem acesso aos dados de sessão e prevenindo vulnerabilidades de roubo de credenciais.
    \item \textbf{Proteção de Rotas via Middleware:} O acesso às páginas privadas (tanto da área reservada de cliente como da administração) é interceptado por um componente de filtragem do Next.js (Middleware). Este componente valida a integridade do token, o seu prazo de validade e o perfil associado ao utilizador antes de permitir a renderização das páginas.
\end{itemize}

\subsubsection{Segurança nas Aplicações .NET (Backoffice e Quiosque)}
As ferramentas administrativas internas e de pesagem utilizam um modelo de autenticação adaptado às condições do ambiente onde são utilizadas:
\begin{itemize}
    \item \textbf{Autenticação por Código de Acesso (PIN):} No terminal quiosque de pesagem física, os operadores autenticam-se através de um código numérico pessoal. O sistema armazena este código encriptado recorrendo a algoritmos de dispersão (\textit{hashing}) seguros, impedindo a sua extração direta em caso de acesso indevido à base de dados estrutural.
    \item \textbf{Autorização de Quiosques (Pairing Code):} Para prevenir acessos não autorizados a partir de dispositivos físicos desconhecidos, a comunicação de novos quiosques com a API exige um processo de emparelhamento estruturado. O administrador gera um código temporário no backoffice e introduz este código no novo terminal para associar permanentemente o identificador do quiosque à lista de dispositivos autorizados.
    \item \textbf{Controlo de Acesso Baseado em Perfis:} A autorização interna baseia-se num sistema de controlo de acessos baseado em permissões granulares. Cada cargo profissional possui um conjunto de permissões ativas estruturado em formato JSON. Após a autenticação, a API central emite um token de sessão interno que transporta estas permissões e que é validado em cada pedido efetuado pelas aplicações Blazor.
\end{itemize}

\subsection{Resultados do Website Institucional}

Nesta secção, apresentam-se os resultados visuais e operacionais obtidos com a implementação do website institucional da empresa, ilustrando o estado final das principais interfaces e fluxos disponíveis.

A página principal do website (Figura~\ref{fig:homepage}) constitui a face pública da empresa. O seu desenho foca-se na facilidade de navegação e na promoção visual da marca, estruturando-se com um cabeçalho fixo com efeito de transparência, secções informativas de serviços e estatísticas gerais sobre o impacto ecológico das atividades de gestão de resíduos da empresa.

\begin{figure}[H]
    \centering
    \framebox[0.9\textwidth]{\parbox{0.85\textwidth}{\centering\vspace{1cm}[Captura de Ecrã: Homepage]\vspace{1cm}}}
    \caption{Interface pública da Página Inicial (Homepage) do website institucional.}
    \label{fig:homepage}
\end{figure}

A interface da área reservada de cliente (Figura~\ref{fig:dashboard}) oferece aos utilizadores autenticados um painel consolidado com a síntese de toda a sua atividade recente. Este ecrã organiza cartões informativos com métricas agregadas e atalhos rápidos para as secções operacionais principais, facilitando a gestão do histórico e pedidos.

\begin{figure}[H]
    \centering
    \framebox[0.9\textwidth]{\parbox{0.85\textwidth}{\centering\vspace{1cm}[Captura de Ecrã: Área de Cliente - Dashboard]\vspace{1cm}}}
    \caption{Ecrã do Painel de Controlo (Dashboard) na área reservada de cliente.}
    \label{fig:dashboard}
\end{figure}

O formulário de submissão de pedidos de recolha (Figura~\ref{fig:novopedido}) permite ao cliente registar uma nova solicitação de forma guiadca. O ecrã dispõe de controlos de seleção de locais de recolha, escolha paginada de materiais regulamentares e uma área dedicada ao carregamento de ficheiros para anexar imagens ou comprovativos de suporte.

\begin{figure}[H]
    \centering
    \framebox[0.9\textwidth]{\parbox{0.85\textwidth}{\centering\vspace{1cm}[Captura de Ecrã: Novo Pedido de Recolha]\vspace{1cm}}}
    \caption{Formulário de submissão de um novo pedido de recolha de resíduos.}
    \label{fig:novopedido}
\end{figure}

Por fim, o painel de administração de notícias (Figura~\ref{fig:adminnews}) é acessível a utilizadores administrativos e disponibiliza uma interface para gerir os conteúdos editoriais do site. O ecrã exibe a tabela de publicações, estatísticas básicas de leitura e botões de ação para abrir o editor de texto estruturado ou o histórico de revisões de artigos.

\begin{figure}[H]
    \centering
    \framebox[0.9\textwidth]{\parbox{0.85\textwidth}{\centering\vspace{1cm}[Captura de Ecrã: Painel Admin - Notícias]\vspace{1cm}}}
    \caption{Painel administrativo de gestão de publicações e notícias do website.}
    \label{fig:adminnews}
\end{figure}

\subsection{Resultados das Apps .NET}

Abaixo são descritos os ecrãs e resultados operacionais desenvolvidos para as aplicações em .NET 10 (Backoffice administrativo e quiosque tátil de balança), que suportam a gestão operacional interna.

O dashboard central da aplicação de Backoffice (Figura~\ref{fig:backoffice}) concentra as listagens de clientes e agendamentos logísticos. A interface exibe tabelas com dados dinâmicos e ordenáveis, contendo menus de ação contextual ativados por toque longo, que facilitam a edição e o processamento de recolhas e a visualização do histórico de auditoria.

\begin{figure}[H]
    \centering
    \framebox[0.9\textwidth]{\parbox{0.85\textwidth}{\centering\vspace{1cm}[Captura de Ecrã: Backoffice - Gestão de Recolhas]\vspace{1cm}}}
    \caption{Interface principal de gestão de recolhas no painel de Backoffice.}
    \label{fig:backoffice}
\end{figure}

O ecrã de entrada do quiosque tátil de balança (Figura~\ref{fig:kiosklogin}) é otimizado para ecrãs táteis em ambiente industrial. Apresenta botões sobredimensionados e um teclado numérico virtual que permite aos operadores introduzir o seu código de acesso pessoal PIN sem necessidade de periféricos físicos, garantindo a proteção e o controlo de acessos.

\begin{figure}[H]
    \centering
    \framebox[0.9\textwidth]{\parbox{0.85\textwidth}{\centering\vspace{1cm}[Captura de Ecrã: Quiosque - Login]\vspace{1cm}}}
    \caption{Interface de login tátil por PIN do quiosque físico de balança.}
    \label{fig:kiosklogin}
\end{figure}

O assistente de registo de pesagens no quiosque (Figura~\ref{fig:kioskweigh}) orienta o funcionário ao longo das fases de pesagem física. O assistente exibe em tempo real o valor lido a partir da balança industrial ligada por porta série, permitindo a seleção do tipo de resíduo e a confirmação do registo estruturado da recolha.

\begin{figure}[H]
    \centering
    \framebox[0.9\textwidth]{\parbox{0.85\textwidth}{\centering\vspace{1cm}[Captura de Ecrã: Quiosque - Pesagem]\vspace{1cm}}}
    \caption{Assistente de pesagem com leitura de peso automática via porta série.}
    \label{fig:kioskweigh}
\end{figure}

\section{Manutenção e Triagem de Hardware}

Paralelamente às tarefas de desenvolvimento de software, foram realizadas diversas atividades práticas e manuais no departamento técnico da empresa, focadas no processamento, triagem e recondicionamento de REEE. Estas atividades visam a economia circular, procurando testar e recuperar componentes funcionais antes de os encaminhar para reciclagem de matérias-primas.

\subsection{Testes e Diagnóstico de Memórias RAM}

A triagem de módulos de memória RAM constitui um passo essencial na recuperação de computadores obsoletos ou danificados que são entregues na empresa. Foram efetuados testes sistemáticos a memórias de tecnologia DDR3 e DDR4, tanto de formato DIMM (computadores de secretária) como SO-DIMM (portáteis). O processo estruturou-se da seguinte forma:
\begin{itemize}
    \item \textbf{Inspeção Visual Inicial e Limpeza:} Verificação de danos físicos no circuito impresso, contactos oxidados ou ausência de componentes passivos (como condensadores SMD). Os contactos dourados foram limpos de forma a remover qualquer oxidação.
    \item \textbf{Verificação de Funcionamento Prático:} Instalação de cada módulo individual numa placa-mãe/computador de teste compatível e ligação do equipamento. Caso o sistema arrancasse com sucesso e apresentasse imagem no ecrã (sinal de vídeo/POST), o módulo era validado como funcional.
    \item \textbf{Classificação e Armazenamento:} Os módulos aprovados foram classificados de acordo com a sua capacidade (GB) e tecnologia (DDR3 ou DDR4) e organizados para reutilização em equipamentos recondicionados. Módulos que não geravam imagem foram descartados e encaminhados para reciclagem de REEE.
\end{itemize}

\subsection{Testes de Unidades de Armazenamento (SATA SSD e M.2)}

Com a crescente substituição de discos mecânicos por unidades de estado sólido, a empresa recebe volumes significativos de dispositivos de armazenamento SSD, tanto com interface SATA de 2,5 polegadas como no formato M.2 (SATA e NVMe). Para avaliar a sua viabilidade de reutilização, seguiu-se o seguinte protocolo:
\begin{itemize}
    \item \textbf{Ligação e Apagamento de Dados:} Ligação das unidades a computadores de teste ou adaptadores dedicados. Efetuou-se uma formatação de baixo nível ou apagamento seguro para garantir a eliminação de dados anteriores e preservar a confidencialidade das informações de antigos utilizadores.
    \item \textbf{Análise de Saúde S.M.A.R.T.:} Execução do software Hard Disk Sentinel para ler os dados de integridade recorrendo à tecnologia SMART do disco. Foram avaliados a percentagem de saúde do dispositivo, o número de setores danificados e o tempo de atividade acumulado.
    \item \textbf{Destino das Unidades:} Unidades em excelente estado de saúde (normalmente com saúde superior a 80%) foram catalogadas e separadas para recondicionamento. Discos inoperacionais ou com níveis severos de degradação foram separados para reciclagem.
\end{itemize}

\subsection{Triagem e Separação de Baterias de Portáteis}

As baterias de iões de lítio (Li-ion) provenientes de computadores portáteis exigem um processo de triagem cuidadoso devido aos riscos de segurança associados ao seu armazenamento e transporte (risco de curto-circuito ou fuga térmica). As atividades realizadas nesta área englobaram:
\begin{itemize}
    \item \textbf{Triagem por Estado Físico:} Inspeção visual minuciosa para deteção de baterias com fissuras, fugas de eletrólito ou sinais de inchaço (deformação física). As baterias danificadas ou inchadas foram imediatamente isoladas e depositadas num bidão fechado e seguro, próprio para resíduos perigosos e inflamáveis.
    \item \textbf{Isolamento dos Terminais:} Aplicação sistemática de fita isoladora sobre os pinos de contacto de cada bateria para evitar o contacto acidental entre os terminais e mitigar riscos de curto-circuito durante o manuseamento e armazenamento.
    \item \textbf{Separação por Marca:} Organização das baterias em bom estado físico de acordo com o fabricante/marca do portátil correspondente (ex: HP, Dell, Asus, Lenovo), facilitando o seu emparelhamento com futuros equipamentos a recondicionar.
\end{itemize}

\subsection{Separação e Validação de Transformadores de Alimentação}

A triagem de transformadores (carregadores de computadores portáteis e periféricos diversos) permitiu recuperar centenas de unidades funcionais, evitando o desperdício de componentes elétricos de alta qualidade. O trabalho manual desenvolvido incluiu:
\begin{itemize}
    \item \textbf{Inspeção de Segurança Elétrica:} Exclusão de transformadores que apresentassem danos físicos visíveis na carcaça plástica ou cabos descarnados, dobrados ou cortados.
    \item \textbf{Teste Dinâmico de Alimentação:} Para verificar a integridade da entrega de energia, utilizou-se um computador de testes que se encontrava sem componentes internos como memórias RAM ou processador. Ao ligar o transformador à entrada de alimentação do computador de testes, a receção de energia fazia com que a placa-mãe emitisse um sinal sonoro (apito). A emissão deste som confirmava que o transformador estava a fornecer energia em conformidade.
    \item \textbf{Separação por Marca e Tipo de Conetor:} Triagem e catalogação das unidades aprovadas segundo a marca e o formato geométrico do conector de entrada (ex: conectores cilíndricos de diferentes diâmetros, fichas retangulares proprietárias ou interfaces USB-C), preparando-as para acompanhar equipamentos recondicionados ou para a reciclagem de cobre no caso dos componentes danificados.
\end{itemize}

\cleardoublepage
\chapter{Sugestões, Comentários e Propostas de Melhoria}\label{chap5}
\thispagestyle{fancy}

Neste capítulo é realizada uma análise reflexiva e crítica sobre o decorrer do estágio curricular, integrando as sugestões organizacionais, comentários sobre a experiência de trabalho, as dificuldades enfrentadas, as limitações identificadas no sistema e as propostas de trabalho futuro.

\section{Sugestões}

Relativamente ao funcionamento do estágio e à sua articulação académica, sugere-se maior flexibilidade na definição inicial de cronogramas para projetos com integração de múltiplos módulos. Projetos que envolvem quiosques para ecrãs táteis e controlo de fluxos logísticos estão sujeitos a adaptações frequentes e reestruturação de bases de dados locais e remotas. Um período de margem oficial para testes de integração em ambiente real diminuiria riscos no planeamento.

\section{Comentários}

A integração na 20 Recolher revelou-se uma experiência importante de aprendizagem profissional. O contacto direto com o setor de gestão de resíduos permitiu compreender a importância da digitalização e da automatização de processos logísticos críticos. A colaboração estreita entre os dois estagiários facilitou a discussão de decisões de arquitetura e fomentou boas práticas de trabalho em equipa, controlo de versões e documentação de APIs.

\section{Dificuldades Encontradas e Soluções Adotadas}

Durante o ciclo de desenvolvimento, a equipa enfrentou três dificuldades principais:
\begin{itemize}
    \item \textbf{Gestão de Estados Concorrentes nos Pedidos:} A existência de múltiplos acessos simultâneos por parte de clientes e operadores gerava inconsistências temporárias no estado das recolhas (por exemplo, dois operadores a tentarem aprovar o mesmo pedido). Para resolver isto, foi implementada uma validação ao nível da API baseada em controlo de concorrência otimista e tratamento de exceções no EF Core.
    \item \textbf{Migração de Arquitetura Monolítica para Distribuída:} Numa fase inicial, o backoffice e o terminal de pesagem foram idealizados numa única aplicação Blazor. Contudo, isso gerava duplicação de acessos à base de dados. Adotou-se como solução isolar o acesso aos dados numa API central baseada em \textit{controllers} RESTful. Esta mudança exigiu a refatoração da comunicação para chamadas HTTP assíncronas em JSON, o que estabilizou o sistema.
    \item \textbf{Gestão de Cache no Supabase e Atualizações Realtime:} Garantir que as notificações de chat e o estado dos pedidos fossem atualizados sem sobrecarregar a base de dados com consultas repetitivas (\textit{polling}) foi complexo. A solução consistiu em configurar o cliente Realtime do Supabase baseado em WebSockets apenas para eventos específicos (mensagens de chat), mantendo o restante fluxo sob pedidos REST padrão.
\end{itemize}

\section{Limitações do Projeto}

Apesar de o ecossistema tecnológico ter atingido os objetivos estipulados, persistem algumas limitações:
\begin{itemize}
    \item \textbf{Dependência de Alojamento Cloud (Supabase/Vercel):} Dado que a base de dados centralizada e o website estão alojados em infraestruturas cloud externas, o ecossistema necessita de ligação permanente à Internet. No caso de uma falha de rede física no armazém da empresa, o quiosque de pesagem não conseguirá registar novas pesagens de forma automática na API.
    \item \textbf{Falta de Módulo de Faturação Direta Integrado:} Atualmente, o sistema extrai listagens em formatos estruturados (CSV/PDF) para permitir a faturação manual. O ecossistema não possui uma integração direta via Web Services com o software de faturação certificado da empresa (ex: Primavera ou Sage).
\end{itemize}

\section{Propostas de Melhoria (Trabalho Futuro)}

Para futuros desenvolvimentos e evolução contínua da plataforma, propõem-se as seguintes melhorias técnicas:
\begin{itemize}
    \item \textbf{Modo Offline com Sincronização Posterior no Quiosque:} Implementar um mecanismo de persistência de pesagens local no terminal (usando bases de dados SQLite locais ou armazenamento local encriptado) que permita continuar a operar sem ligação à Internet e envie os dados acumulados de forma assíncrona assim que a ligação à API for restabelecida.
    \item \textbf{Integração de APIs de Software de Faturação Certificado:} Desenvolver conectores dedicados para comunicar com o ERP da empresa, automatizando a emissão imediata da guia de transporte e da fatura após a validação da pesagem no terminal quiosque.
    \item \textbf{Notificações Push e Aplicação Móvel para Clientes:} Expandir a presença web do cliente para uma aplicação híbrida móvel ou implementar notificações push nativas no navegador, alertando instantaneamente os utilizadores quando uma recolha for efetuada ou quando for recebida resposta do suporte técnico.
\end{itemize}

\cleardoublepage
\chapter{Conclusões}\label{chap6}


\section{Balanço do Estágio}
% Descreve a apreciação geral do estágio, relacionando-a com as tuas expetativas iniciais.
% Explica de que modo esta experiência contribuiu para a tua valorização profissional e pessoal.

% A desenvolver

\section{Objetivos Atingidos}
% Confirmação dos objetivos alcançados e análise dos desvios ocorridos em relação ao plano inicial.

% A desenvolver
\cleardoublepage
\chapter{Conclusão}
\label{chap7}
\thispagestyle{fancy}
\acresetall


%%%  %%%%%%%%%%%%%%%%%%%%



%%%%%%%%%%%%%%%%%%%%%%%%%%%%%%%%%%%%%%%%%%%%%%%%%%%%%%%%%%%%%%%%%%%%%
%           BIBLIOGRAFIA

\nocite{*}
% --- FIM DA BIBLIOGRAFIA ---

%\printbibliography
\cleardoublepage
\printbibliography[heading=bibintoc,title={Referências}]

%%%%%%%%%%%%%%%%%%%%%%%%%%%%%%%%%%%%%%%%%%%%%%%%%%%%%%%%%%%%%%%%%%%%%
%                               ANEXOS

\appendix
\chapter{Excertos de Código e Configurações}\label{ap:codigo}
\thispagestyle{fancy}

Neste anexo, apresentam-se alguns excertos de código e configurações críticas desenvolvidas para o ecossistema tecnológico, evidenciando as decisões de integração e segurança tomadas durante o projeto.

\section{Configuração do Proxy Reversível em Next.js}\label{ap:proxy}

De forma a mascarar a localização real da API central e permitir a transmissão segura de cookies com a diretiva \textit{HttpOnly}, configurou-se uma regra de reencaminhamento no ficheiro de configuração do Next.js. O trecho abaixo ilustra como os pedidos direcionados para \texttt{/api/v1/*} são mapeados internamente para a API em .NET:

\begin{lstlisting}[language=HTML, caption={Configuração de rewrites de rede em Next.js.}]
// next.config.mjs
const nextConfig = {
  async rewrites() {
    return [
      {
        source: '/api/v1/:path*',
        destination: process.env.NEXT_PUBLIC_API_URL + '/api/v1/:path*',
      },
    ];
  },
};

export default nextConfig;
\end{lstlisting}

\section{Configuração de Políticas CORS no .NET API}\label{ap:cors}

Para garantir que apenas os clientes autorizados conseguem obter dados da API, implementou-se uma política de CORS estrita na inicialização do serviço .NET, limitando o acesso ao domínio do portal cliente:

\begin{lstlisting}[language=HTML, caption={Configuração de políticas CORS na API em C\#.}]
// Program.cs
var builder = WebApplication.CreateBuilder(args);

builder.Services.AddCors(options =>
{
    options.AddPolicy("PortalClientePolicy", policy =>
    {
        policy.WithOrigins(builder.Configuration["AllowedOrigins:PortalCliente"] 
                ?? "https://portal.20recolher.pt")
              .AllowAnyHeader()
              .AllowAnyMethod()
              .AllowCredentials(); // Permitir envio de tokens de cookies
    });
});

var app = builder.Build();
app.UseCors("PortalClientePolicy");
\end{lstlisting}

\section{Mecanismo de Filtro Global de Soft-Delete no EF Core}\label{ap:softdelete}

Para automatizar a exclusão lógica ao nível da base de dados sem necessidade de adicionar cláusulas condicionais repetitivas nas consultas do Entity Framework Core, configurou-se um filtro de consulta global no contexto de dados:

\begin{lstlisting}[language=HTML, caption={Configuração de filtro global de consulta para soft-delete no EF Core.}]
// ApplicationDbContext.cs
protected override void OnModelCreating(ModelBuilder modelBuilder)
{
    base.OnModelCreating(modelBuilder);

    // Aplicar filtro de eliminação lógica para todas as entidades que herdam de ISoftDeletable
    foreach (var entityType in modelBuilder.Model.GetEntityTypes())
    {
        if (typeof(ISoftDeletable).IsAssignableFrom(entityType.ClrType))
        {
            modelBuilder.Entity(entityType.ClrType)
                .HasQueryFilter(ConvertFilterExpression(entityType.ClrType));
        }
    }
}
\end{lstlisting}

\chapter{Registo Fotográfico do Lab de Hardware}\label{ap:fotos}
\thispagestyle{fancy}

Neste anexo, apresentam-se fotografias do espaço de trabalho e dos processos manuais de diagnóstico e manutenção efetuados no departamento de hardware da empresa.

\begin{figure}[H]
    \centering
    \confidencial % Indica que as fotografias são de cariz operacional privado
    \vspace{2cm}
    \textbf{[FOTOGRAFIA 1: Área de Triage e Diagnóstico de REEE]} \\
    \textit{Espaço reservado aos testes de RAM, discos e transformadores.}
    \vspace{2cm}
    \caption{Bancada de ensaio e computadores de teste utilizados no laboratório.}
    \label{fig:foto_bancada}
\end{figure}

\begin{figure}[H]
    \centering
    \confidencial
    \vspace{2cm}
    \textbf{[FOTOGRAFIA 2: Processo de Recondicionamento de Equipamentos]} \\
    \textit{Etapa de higienização interna e aplicação de composto térmico na CPU.}
    \vspace{2cm}
    \caption{Manutenção física e preparação estética de portátil recondicionado.}
    \label{fig:foto_manutencao}
\end{figure}


%%%%%%%%%%%%%%%%%%%%%%%%%%%%%%%%%%%%%%%%%%%%%%%%%%%%%%%%%%%%%%%%%%%%%
%                       ÚLTIMA PÁGINA

\blankpage{}

\justifying    

O presente documento constituído por \pageref{LastPage} páginas numeradas, com anexos, foi elaborado pelo signatário no âmbito da Formação em Contexto de Trabalho (Estágio) do Curso Técnico Superior Profissional (CTeSP) em Tecnologias e Programação de Sistemas de Informação.\\

\begin{center}

Entregue em \_\_ de \_\_\_\_\_\_\_\_\_\_\_\_\_\_\_\_ de 2026

\vspace{1.5cm}

Nome do aluno: Guilherme Laranjeiro Ventura

\vspace{1.5cm}

Assinatura (Obrigatória)

\end{center}

\blankpage{}

\end{document}